\chapter{Introduction}

This book is a guided tour on how to use \textsf{WP} plug-in of
\textsf{Frama-C} for proving \textsf{C} programs annotated with
\textsf{ACSL} notations. It is based on the excellent ``ACSL By
Example'' book produced by the Fraunhofer FIRST Institute for 
the Device-Soft project\footnote{Version 7.1.0 of December 2011, 
  see \url{http://www.first.fraunhofer.de}}.

We assume the reader to be familiar with \textsf{ACSL} in general and
already equipped with the \textsf{WP} plug-in, which is distributed
with \textsf{Frama-C} releases. Please refer to the \textsf{WP} user's
manual\footnote{\url{http://frama-c.com/download/frama-c-wp-manual.pdf}}
for installation and general overview of the plug-in.

\section{Library}

The case studies presented in this document are exact copies of the
ones presented in the original ``ACSL By Example'' book. Sometimes, we
indicate some modifications of the specification that makes
\textsf{WP} better prove the programs.

All examples uses the following \emph{library} of \textsf{C} and
\textsf{ACSL} definitions:

\listingname{library.h}
\cinput{library.h}

\section{Examples}

Source of case studies are generally presented in two separated files:
one for the header and specification of the algorithm, and one of its
implementation.  To ease the presentation in the book, we have omitted
the necessary \texttt{include} lines from sources.

Thus, a header file \texttt{A.h} should be enclosed by the following lines:
\begin{ccode}
  #ifndef _A_H
  #define _A_H
  #include "../library.h"
  [...]
  #endif
\end{ccode}

Similarly, an implementation file \texttt{A.c} should start by including its header:
\begin{ccode}
  #include "A.h"
  [...]
\end{ccode}
